% Laget automatisk av elevark/lag_elevark.py fra sim/rekyl.html. Endringer her blir overskrevet.
\documentclass[11pt]{article}
\usepackage{elevark}
\begin{document}
\elevarkhode{Fysikk 1 · Bevegelsesmengde og støt}{Rekyl}{https://lektorbo.no/sim/rekyl.html}{Simuleringen}

To skøyteløpere står i ro på isen og dytter fra hverandre. Hvem får størst fart, og hvem glir lengst? Det samme skjer når du kaster en tung ball fra skøyter, og når en kanon skyter.

\section*{Oppgaver}

Før dyttet står alt i ro, så den samlede bevegelsesmengden er null. Hva betyr det for bevegelsesmengden etterpå?

\begin{oppgaveliste}
\oppgave \fase{gjett}{Gjett}En voksen på 80~kg og et barn på 30~kg dytter fra hverandre. Hvem dytter hardest på den andre? Hvem får størst fart?
\svar
\oppgave \fase{test}{Test}Kjør simuleringen. Sammenlign fartene og bevegelsesmengdene. Hva er den samlede bevegelsesmengden etter dyttet?
\svar
\oppgave \fase{test}{Regn}De dytter med 200~N i 0,40~s. Regn ut impulsen og farten til hver av dem. Hvor langt glir hver av dem når friksjonstallet er 0,015?
\svar
\oppgave \fase{forklar}{Forklar}Barnet glir mer enn sju ganger så langt som den voksne, selv om farten bare er 2,7 ganger så stor. Forklar hvorfor.
\svar
\oppgave \fase{test}{Undersøk}Velg «Kast fra skøyter». Hvordan endrer farten din seg når du kaster en tyngre ball, eller kaster fortere? Hvorfor er det vanskelig å komme seg av en helt glatt is uten å ha noe å kaste?
\svar
\oppgave \fase{forklar}{Forklar}Velg «Kanon på hjul». Kanonen og kula får like stor bevegelsesmengde. Hvorfor får kula likevel nesten all den kinetiske energien?
\svar
\oppgave \fase{utfordring}{Utfordring}En rakett i verdensrommet har ingenting å skyve mot. Bruk det du har sett, til å forklare hvordan den likevel kan øke farten.
\svar
\end{oppgaveliste}

\end{document}
